\documentclass[12pts]{article}

%Importamos los paquetres
\usepackage[utf8]{inputenc} %Configura caracteres especiales
\usepackage[spanish]{babel} %Configura en español los valores predeterminados
\usepackage{amsmath} %Para manejo de formulas
\usepackage{array} %Mejor manejo de las tablas
\usepackage{geometry} %Margenes
\geometry{margin=2.5cm} %Modifica el margen de la hoja

\newcommand{\minombre}{Yessica Nahomi}
\newcommand{\nombrepractica}{Practica 4: Tablas, entornos y comandos}
\newcommand{\universidad}{UASLP-FEPZM}

\newenvironment{ejercicio}
{
\vspace{0.3cm}
\noindent\textbf{Ejercicio:}
\begin{itshape}%Modifica a italica el texto del ejercicio
}
{
\end{itshape}
\vspace{0.3}
}
\begin{document}
\begin{center}
{\LARGE \nombrepractica\\[5.5cm]
{\LARGE \minombre}\\[5.5cm]
{\LARGE \universidad}\\[5.5cm]
{\LARGE \today}
\end{center}


\section{Introduccion}
En esta practica aprenderemos a crear tablas y usar \textbf{environments}, incluyendo environments personalizados.

\section{Uso de tablas}
\begin{table}[h]
\centering
\begin{tabular}{c|c|c}
\hline
Alumno&Materuia&Calificacion\\
Pedro&Matematicas&8.0\\
Maria&Español&10.0\\


\end{tabular}
\caption{Calificaciones de alumnos}
\label{tab:calificaciones}
\end{table}

\subsection{Tablas con columnas combinadas}
\begin{table}[h]
    \centering
    \begin{tabular}{c|c|c}
    \hline
    \multicolumn{3}{|c|}{Horario}\\ \hline
    Dia&Horario&Materia \\ \hline
    Lunes&10 a 11&Matematicas \\ \hline
    Martes&12 a 13&Fisica \\ \hline

    \end{tabular}
    \caption{Horario semanal}
    \label{tab:placeholder}
\end{table}


\section{Environment personalizado}
Ahora usaremos nuestro environment \textbf{ejercicio}

\begin{ejercicio}
    Crear una tabla con al menos 4 filas y 3 columnoas, rellenala con contenido afin a un tema en especifico.

\end{ejercicio}

\section*{Tabla sobre Lenguajes de Programación}

\begin{center}
\begin{tabular}{|c|c|c|}
\hline
\textbf{Lenguaje} & \textbf{Uso Principal} & \textbf{Año de Creación} \\
\hline
Python & Inteligencia Artificial & 1991 \\
\hline
Java & Aplicaciones Empresariales & 1995 \\
\hline
C++ & Desarrollo de Videojuegos & 1985 \\
\hline
JavaScript & Desarrollo Web & 1995 \\
\hline
\end{tabular}
\end{center}



\end{document}